\section{Chain-of-Thought Signatures}
\label{sec:cot}

Section~\ref{sec:compare} evaluated \emph{whether} an intervention produces a selective safety failure; we now ask \emph{how} that failure appears in visible reasoning. This distinction matters for a safety-mechanism claim. Two interventions may produce similar harmful answers while leaving different trace signatures: one may reduce explicit harm recognition or refusal reasoning, whereas another may preserve these signals but increase downstream execution. We do not treat the latter as normatively preferable; rather, the distinction reveals whether answer-level failure is accompanied by less visible safety reasoning or by compliance despite retained safety signals. We analyze this using Monitorability Retention, $\Sreason$ (\S\ref{sec:axes}), a per-sentence safety-reasoning judge, and a $12$-label pathway judge (\S\ref{sec:judges}). The section makes three points: (i) the interventions do not systematically reduce trace monitorability (\S\ref{subsec:cot-monitor}), (ii) harmful completions often retain visible safety signals while shifting toward downstream execution (\S\ref{subsec:cot-pathway}), and (iii) this signature is clearest and scales with steering dose on explicit-trace models (\S\ref{sec:cot-reasoning}). We report several additional observations supported by the pathway analysis in appendix \ref{sec:appendix-cot-other}.

\subsection{Monitorability Gap as a Measure}
\label{subsec:cot-monitor}

Following the paired trace-answer protocol in \S\ref{sec:cot-monitoring}, we ask whether an intervention produces harmful answers from benign-looking
traces. We report the signed diagnostic $\gapm=U-O=\mathrm{asr}_{\text{final}}-\mathrm{asr}_{\text{cot-pred}}$ alongside $\Sreason$. A positive gap indicates that covert failures outnumber over-warnings, whereas a zero or negative gap indicates that the trace monitor is at least as likely to over-warn as to miss harmful output. This signed gap is a descriptive net measure of under- versus over-warning on the shared paired prompt set, not a prompt-level agreement score: because $U$ and $O$ can cancel, we also inspect paired agreement $A$ and trace FNR when interpreting individual cells.

We find no systematic shift toward covert failure. Across all $100$ intervention cells, including the early-response-prefix sensitivity arms, the
gap is non-positive in $96$ cells. Only $4/100$ cells have a positive gap, with a maximum of $+0.036$. The mean gap by family ranges from $-0.107$ for safety head ablation to $-0.075$ for directional ablation, and the most negative cells are safety head ablations on \model{Llama-3.1-8B}, where the gap reaches $-0.448$. These are the same low-Potency cells identified in \S\ref{sec:compare}: answer coherence degrades, while the trace or early-response prefix continues to flag the request as harmful. We therefore find no systematic evidence that higher Potency comes at the cost of less revealing traces (Figure~\ref{fig:gap-family}).

\begin{figure}[t]
  \centering
  \includegraphics[width=\linewidth]{figures/fig_gap_by_family.png}
  \caption{\textbf{No systematic shift toward covert failure.} Distribution of the signed monitorability gap $\gapm$ by family across $100$ intervention
  cells (five model/configuration arms $\times$ two datasets $\times$ ten non-baseline conditions). Negative values mean that over-warnings outnumber
  covert failures at the cell level; they do not imply that every individual trace is classified correctly. The negative tail is driven by low-Potency,
  coherence-degrading head-ablation cells on \model{Llama-3.1-8B}.}
  \label{fig:gap-family}
\end{figure}

Monitorability also varies independently of how much harm an intervention causes. Across cells, the gap is uncorrelated with Potency (Spearman
$\rho\!\approx\!0.00$). Within the two explicit-trace models, \model{Qwen3-8B} and \model{OLMo-3-Think}, the association is mildly negative
($\rho=-0.28$ to $-0.33$). The monitorability diagnostics are therefore not redundant with $\Ppot$.

\subsection{Pathway Signatures Shift Toward Downstream Execution}
\label{subsec:cot-pathway}

The pathway judge provides a finer-grained view of the visible reasoning traces. We report a compact signature of which reasoning stages change under
each family (Table~\ref{tab:cot-signature}), summarizing five indicators that span harm recognition, refusal dynamics, and execution. The full $12$-label
matrix is reported in Appendix~\ref{sec:appendix-pathway}.

\begin{table}[t]
  \centering \small
  \setlength{\tabcolsep}{4pt}
  \resizebox{\columnwidth}{!}{%
  \begin{tabular}{@{}lccccc@{}}
    \toprule
    \textbf{Family} & recog. & ref.\ init. & ref.\ supp. & exec. & op.\ detail \\
    \midrule
    Steering              & 0.52 & 0.76 & 0.11 & 0.35 & 0.27 \\
    Directional ablation  & 0.53 & 0.91 & 0.05 & 0.31 & 0.20 \\
    Neuron ablation       & 0.60 & 0.95 & 0.09 & 0.24 & 0.15 \\
    SHIPS (heads)         & 0.60 & 0.93 & 0.10 & 0.24 & 0.14 \\
    \bottomrule
  \end{tabular}}
  \caption{\textbf{Pathway signature by family on the explicit-trace models.} Values pool harmful-among-clean completions from \model{OLMo-3-Think} and
    \model{Qwen3-8B}. Recognition and refusal initiation remain common, while explicit refusal suppression is rare. Harmful completions are therefore
    more often associated with downstream execution and operational detail than with explicit suppression of an initial refusal.}
  \label{tab:cot-signature}
\end{table}

Two results are consistent across intervention families. First, explicit refusal suppression remains rare ($0.05$--$0.11$), while recognition
($0.52$--$0.60$) and refusal initiation ($0.76$--$0.95$) remain common. Harmful completions, therefore, often retain visible recognition and an initial
refusal before proceeding to harmful content. Second, downstream labels track induced harm: operational detail and execution correlate positively with
Potency (Spearman $\rho=+0.49$ and $+0.39$), whereas recognition and refusal initiation correlate negatively ($\rho=-0.42$ and $-0.37$); refusal suppression is only weakly related ($\rho=-0.14$). Relative to the non-intervention baseline, the interventions change refusal suppression by at most $0.01$, while shifting recognition downward and execution upward. This pattern is more consistent with downstream bypass in visible reasoning than with a wholesale loss of harm recognition. It describes the process signature in the trace rather than establishing the location of the underlying latent
computation.

The radar in Figure~\ref{fig:radar} complements the pooled result with the strongest condition from each family on \model{OLMo-3-Think}. Directional
ablation, SHIPS, and neuron ablation remain comparatively close to baseline, whereas the strongest steering condition shifts sharply away from recognition
and refusal initiation toward procedural planning and operational detail.

\begin{figure}[t]
  \centering
  \includegraphics[width=\linewidth]{figures/pathway_radar_overlay_olmo3_7b_think.png}
  \caption{\textbf{Strong steering produces the largest pathway shift on \model{OLMo-3-Think}.} The figure compares the baseline with the strongest
  condition from each family across the $12$ pathway labels (dataset mean over parsed CoT prefixes). Directional ablation, SHIPS, and neuron ablation remain close to the baseline profile. Steering at $\alpha=1.5$ reduces risk recognition, policy awareness, and refusal initiation while increasing
  procedural planning and operational detail. This strongest-condition view complements the family averages in Table~\ref{tab:cot-signature}; the full
  per-cell label table is in Appendix~\ref{sec:appendix-pathway}.}
  \label{fig:radar}
\end{figure}

\subsection{The Signature Is Clearest in Explicit-Trace Models}
\label{sec:cot-reasoning}

The pathway interpretation depends on the availability of explicit, parseable reasoning traces. Scoring traces with the per-sentence safety-reasoning judge (\S\ref{sec:judges}), the two explicit-trace models (\model{OLMo-3-Think} and \model{Qwen3-8B}) verbalize safety reasoning in $86$--$90\%$ of traces, compared with $15$--$35\%$ of the early-response prefixes used for the remaining arms. The recognize-but-proceed signature is therefore clearest on the explicit-trace models: recognition and refusal initiation remain visible even when the final answer is harmful. Because the other arms expose only response prefixes, we treat their corresponding analysis as a sensitivity check rather than as private chain-of-thought.

We conduct experiments on the effect of intervention-dosage on trace, pooled across the two reasoning-models. This shows that the increased Potency in steering is accompanied by an intensified pathway-signature, while trace does not change with dosage in component-level intervention. Details can be seen in appendix \ref{sec:appendix:dosage}.
